A radium-226 nucleus emits alpha particles ($m_\alpha = \maO$) with a kinetic energy of \EaO. Inside the nucleus, the alpha particle is held by the strong force; at the edge of the nucleus, the electric repulsion forms a wall about \EbO\ high. The half-life of radium-226 is 1600~years; that of polonium-212, whose alpha particles have \EpO, is \SI{0.3}{\micro\s}.

\begin{abcliste}
\abc Calculate the speed of the alpha particles.

\abc Calculate their de Broglie wavelength.

\abc Could an alpha particle of \EaO\ leave the nucleus according to classical physics? Explain.

\abc Why is the half-life of polonium-212 so much shorter?
\end{abcliste}
